\subsection{A Brief History of Reinforcement Learning}

The reinforcement learning (RL) journey began with early computational work on game engines. \citet{Shannon1950} created chess evaluation functions, anticipating programs that could ``learn from mistakes." \citet{Minsky1954}'s thesis connected neural networks to Skinnerian conditioning, introducing a "second memory" for outcome prediction—foreshadowing value functions. \citet{Samuel1959}'s groundbreaking checkers program implemented rote learning and parameter adjustment based on evaluation differences, presaging temporal difference learning. Concurrently, \citet{Bellman1957} and \citet{Howard1960} developed dynamic programming for sequential decision-making, establishing the mathematical foundation of optimality through backward induction in fully-specified environments. These early contributions identified the core challenges that would drive RL research: credit assignment (determining which actions lead to rewards), exploration (trying new actions while exploiting successful ones), and function approximation (generalizing across state spaces).

Reinforcement learning crystallized as a coherent field through Sutton and Barto's unifying work, which addressed the critical problem of learning without complete environmental models. Their temporal difference (TD) learning method \citep{Sutton1988} solved the credit assignment problem by bootstrapping—updating value estimates based on subsequent predictions rather than waiting for final outcomes—a middle ground between Monte Carlo approaches and dynamic programming. This innovation enabled agents to learn incrementally from partial experiences, a fundamental capability for online learning. Building on this foundation, they synthesized diverse approaches (Q-learning, SARSA, actor-critic methods) under the framework of generalized policy iteration (GPI), where policy evaluation and improvement operate symbiotically \citep{SuttonBarto1998}. 

A significant early framework was the actor-critic architecture, introduced by \citet{BartoSuttonAnderson1983}, which combined a policy representation (actor) with a separate value function (critic). This approach allowed agents to handle both discrete and continuous action spaces while providing a natural implementation of policy gradient methods. \citet{watkins1992_paper}'s Q-learning, with its convergence guarantees, allowed an agent to learn optimal policies despite following a different behavioral policy. This "off-policy" learning allowed the agent to re-use past experiences. Meanwhile, policy gradient methods like REINFORCE \citep{Williams1992} directly optimized policy parameters, being particularly valuable for continuous or high-dimensional action spaces where value-based methods struggled.

The multi-armed bandit problem, a simplified setting where agents make one-shot decisions under uncertainty, provided a formal framework for analyzing the exploration-exploitation trade-off. Lai \& Robbins (1985) established regret lower bounds which was integrated into robust algorithms like Upper Confidence Bound (UCB) by \citet{Auer2002} which had finite-time regret guarantees. 

Function approximation represented the next frontier, as traditional tabular methods couldn't scale to problems with large state spaces. \citet{Tesauro1994}'s TD-Gammon marked a breakthrough: a backgammon-playing program that achieved near-human-expert performance through self-play, using TD learning with neural networks as function approximators. This success demonstrated that RL could scale to domains with proper generalization mechanisms. However, function approximation introduced new stability challenges—the ``deadly triad" of bootstrapping, off-policy learning, and function approximation could cause divergence and instability in learning.

Before the deep learning era, the RL toolkit expanded in different directions. \citet{Sutton1990}'s Dyna architecture unified learning and planning by integrating a learned model with direct reinforcement learning, establishing the model-based RL paradigm that would later influence systems like AlphaGo. Hierarchical frameworks tackled the complexity of long-horizon problems: \citet{Sutton1999Options}'s Options framework formalized temporal abstraction, allowing agents to reason at multiple time scales. \citet{brafman2002rmax}'s R-MAX algorithm provided theoretical guarantees for efficient exploration through the principle of "optimism in the face of uncertainty," automatically balancing exploration and exploitation. Policy optimization advanced through \citet{Kakade2001}'s Natural Policy Gradient, which improved learning efficiency by accounting for the concavity of the loss landscape through Fisher information. \citet{NgRussell2000} formalized Inverse Reinforcement Learning for inferring reward functions from demonstrations, enabling preference learning from expert behavior—a subfield with exactly the same goals as structural econometrics. \citet{BertsekasTsitsiklis1996} bridged RL with control theory through their "Neuro-Dynamic Programming" framework, providing rigorous theoretical analysis and applications to operations research problems.

The deep RL revolution overcame the problem of function approximation and the "deadly triad" that had limited previous scaling efforts. \citet{Mnih2015}'s Deep Q-Networks (DQN) achieved human-level performance on Atari games by combining Q-learning with convolutional neural networks, learning directly from raw pixels without hand-crafted features. Critical innovations that stabilized this integration included experience replay (breaking correlations between consecutive samples) and target networks (reducing update instability). These techniques addressed key problems in applying RL to high-dimensional sensory inputs, enabling end-to-end learning from raw observations to actions. The breakthrough opened a floodgate of innovations, with \citet{Schulman2015}'s Trust Region Policy Optimization and \citet{Schulman2017}'s Proximal Policy Optimization providing similarly stable and reliable policy learning methods.

Game-playing achievements demonstrated how combining deep RL with planning could solve problems of unprecedented complexity. \citet{Silver2016}'s AlphaGo defeated world champion Lee Sedol by integrating supervised learning from human games with reinforcement learning from self-play, marking a milestone that many experts believed was decades away. More remarkably, AlphaGo Zero \citep{Silver2017} achieved superhuman performance without human data, learning entirely through self-play. \citet{Igami2020} noted the economic parallels, interpreting AlphaGo's architecture as combining a supervised learning policy network (similar to Hotz-Miller's conditional choice probability estimator) with a reinforcement learning value network (analogous to conditional choice simulation methods). 

The application of reinforcement learning to language models represents one of its most significant recent developments. Reinforcement Learning from Human Feedback (RLHF) emerged as the dominant paradigm for aligning large language models with human preferences, addressing the challenge of optimizing for implicit human values rather than easily-specified reward functions. Key implementations include OpenAI's InstructGPT \citep{Ouyang2022}, which showed a significantly smaller RLHF-tuned model outperforming much larger base models on human preference metrics, and Anthropic's Constitutional AI \citep{Bai2022}, which developed RL from AI Feedback to reduce the need for human evaluations. These applications demonstrate how an RL framework enables AI systems to better align with human intentions.

\subsection{A Brief History of Structural Econometrics}

Structural econometrics emerged from the Cowles Commission's groundbreaking work in the 1940s, which established the foundation for relating economic theory to statistical analysis. \citet{Haavelmo1944} introduced the probabilistic framework that distinguished between structural and reduced-form parameters, formalizing the identification problem in simultaneous equations. Koopmans developed order and rank conditions determining which parameters are identifiable given model specifications, while Marschak emphasized the importance of "autonomy" in structural equations—relationships that remain invariant under policy interventions. \citet{HoodKoopmans1953} consolidated these advances into a coherent set of tools, introducing estimation techniques like limited and full-information maximum likelihood. This methodological fusion of economic theory with statistical inference became the defining characteristic of structural econometrics.

The field advanced significantly through \citet{McFadden1974}'s discrete choice framework, which formalized the connection between utility maximization theory and probabilistic choice behavior. \citet{Heckman1979}'s sample selection model addressed estimation bias from non-random sample selection, treating participation decisions as part of the economic model. These static frameworks laid groundwork for structural estimation in labor economics, industrial organization, and public finance. By the early 1980s, researchers began transitioning from static to dynamic sequential decision problems. \citet{Wolpin1984} modeled fertility decisions over time, and \citet{Miller1984} developed dynamic job matching models.

A milestone was \citet{Rust1987}'s nested fixed-point (NFXP) algorithm, which formulated and estimated a fully specified Markov decision process for engine replacement decisions. For each trial parameter vector, Rust solved the Bellman equation via fixed-point iteration, then maximized the likelihood of observed choices. This methodological breakthrough showed that full-solution estimation of dynamic discrete choice models was computationally feasible, inspiring subsequent innovations to address the "curse of dimensionality." \citet{BerryLevinsohnPakes1995} revolutionized demand estimation with their random coefficients logit model for differentiated products, inverting market share data to recover underlying consumer utilities while allowing flexible substitution patterns with endogenous prices. This "BLP" approach became the workhorse for static demand estimation in industrial organization, later enhanced through techniques integrating micro-data \citep{Petrin2002} and methodological refinements addressing numerical performance \citep{DubeFoxSu2012}.

The computational burden of repeatedly solving dynamic games spurred methodological innovations in the 2000s. \citet{HotzMiller1993} introduced the conditional choice probability (CCP) method, showing that observed choice probabilities could be inverted to recover value function differences without repeatedly solving the dynamic programming problem. \citet{BajariBenkardLevin2007} adapted this insight to dynamic games through a two-step approach—first estimating policy functions from data, then finding structural parameters that rationalize observed behavior. \citet{AguirregabiriaMira2007} proposed nested pseudo-likelihood estimation, iteratively refining two-step estimates to reduce bias. \citet{PesendorferSchmidtDengler2008} developed value function difference methods using equilibrium conditions directly.

Recent advances have further expanded structural econometrics' capabilities. \citet{WeintraubBenkardVanRoy2008} introduced oblivious equilibrium concepts to handle high-dimensional state spaces with many agents. \citet{CilibertoTamer2009} developed moment inequality methods addressing equilibrium multiplicity through partial identification. Advances in production function estimation \citep{AckerbergCavesFrazer2015} refined proxy methods for handling endogeneity. The 2010s witnessed integration with machine learning techniques, using neural networks as function approximators for value and policy functions. \citet{FernandezVillaverdeNunoPerla2024} demonstrated how deep learning can mitigate the curse of dimensionality in solving high-dimensional economic models. \citet{MaliarMaliarWinant2021} used neural networks trained via stochastic gradient descent for solving dynamic programming problems, while \citet{AtashbarShi2023} applied deep reinforcement learning to macroeconomic models.